\chapter{Infrastructure Layer}
\label{ch:infrastructure}

\textit{\sysname{} runs on three databases and a cache, each chosen for a
specific job.  This chapter explains what they are, how they are wired together
via Docker Compose, how to persist (or destroy) your data, and how Alembic
keeps PostgreSQL schemas in sync across contributors.}


% ─────────────────────────────────────────────────────────────────────
\section{Three Databases}
\label{sec:three-databases}

Most web applications use a single database.  \sysname{} uses three---not out of
complexity fetishism, but because relational data, graph data, and ephemeral
cache data have fundamentally different access patterns.
Table~\ref{tab:database-overview} summarizes the choices.

\begin{table}[ht]
\centering
\caption{Infrastructure services and why each was chosen.}
\label{tab:database-overview}
\begin{tabularx}{\textwidth}{l c X l}
\toprule
\textbf{Service} & \textbf{Port} & \textbf{Purpose} & \textbf{Why Chosen} \\
\midrule
PostgreSQL~18 & 5432 &
    Relational storage for users, projects, decision metadata.
    Schema managed by Alembic migrations. &
    ACID compliance, concurrent access, rich ecosystem.
    Chosen over SQLite for multi-user concurrency. \\
\addlinespace
Neo4j~2025.01 & 7474 / 7687 &
    Knowledge graph: entity nodes, decision traces, relationship edges.
    Queried via Cypher. &
    Native labeled property graph model.
    Cypher makes multi-hop traversals readable. \\
\addlinespace
Redis~7.4 & 6379 &
    Caching (3 tiers with different TTLs), rate limiting,
    session ephemera. &
    Sub-millisecond reads.  Append-only file persistence for
    crash recovery without full durability overhead. \\
\bottomrule
\end{tabularx}
\end{table}

\begin{keyinsight}
PostgreSQL is the system of record for \emph{who} and \emph{what}.  Neo4j is the
system of record for \emph{how things relate}.  Redis is the accelerator---if you
lose its data, nothing is permanently lost; the caches rebuild themselves.
\end{keyinsight}


% ─────────────────────────────────────────────────────────────────────
\section{Docker Compose}
\label{sec:docker-compose}

All three services are defined in a single
\codefile{docker-compose.yml} at the repository root.  Here is the full file:

\begin{lstlisting}[style=yaml, caption={docker-compose.yml (complete).}]
services:
  postgres:
    image: postgres:18-alpine
    container_name: continuum-postgres
    environment:
      POSTGRES_USER: ${POSTGRES_USER:?POSTGRES_USER must be set}
      POSTGRES_PASSWORD: ${POSTGRES_PASSWORD:?POSTGRES_PASSWORD must be set}
      POSTGRES_DB: ${POSTGRES_DB:-continuum}
    ports:
      - "127.0.0.1:5432:5432"
    volumes:
      - postgres_data:/var/lib/postgresql/data
    healthcheck:
      test: ["CMD-SHELL",
             "pg_isready -U $${POSTGRES_USER} -d $${POSTGRES_DB:-continuum}"]
      interval: 10s
      timeout: 5s
      retries: 5
      start_period: 10s
    restart: unless-stopped

  neo4j:
    image: neo4j:2025.01-community
    container_name: continuum-neo4j
    environment:
      NEO4J_AUTH: ${NEO4J_USER:?}/${NEO4J_PASSWORD:?}
      NEO4J_PLUGINS: '["apoc"]'
      NEO4J_dbms_security_procedures_unrestricted: apoc.*
      NEO4J_dbms_memory_heap_initial__size: 512m
      NEO4J_dbms_memory_heap_max__size: 1G
    ports:
      - "127.0.0.1:7474:7474"   # HTTP browser
      - "127.0.0.1:7687:7687"   # Bolt protocol
    volumes:
      - neo4j_data:/data
      - neo4j_logs:/logs
    healthcheck:
      test: ["CMD", "wget", "-q", "--spider",
             "http://localhost:7474"]
      interval: 10s
      timeout: 5s
      retries: 5
      start_period: 30s
    restart: unless-stopped

  redis:
    image: redis:7.4-alpine
    container_name: continuum-redis
    ports:
      - "127.0.0.1:6379:6379"
    volumes:
      - redis_data:/data
    healthcheck:
      test: ["CMD", "redis-cli", "-a",
             "$${REDIS_PASSWORD}", "ping"]
      interval: 10s
      timeout: 5s
      retries: 5
    command: >-
      redis-server --appendonly yes
      --requirepass ${REDIS_PASSWORD:?REDIS_PASSWORD must be set}
    restart: unless-stopped

volumes:
  postgres_data:
  neo4j_data:
  neo4j_logs:
  redis_data:
\end{lstlisting}

A few things to notice:

\begin{itemize}
    \item Every port mapping begins with \ic{127.0.0.1}---the service is only
          reachable from your local machine, never from the network.
    \item Environment variables ending with \ic{:?} are \emph{required}---Docker
          Compose will refuse to start if they are unset.
    \item Neo4j gets a 512\,MB initial heap and a 1\,GB maximum.  The
          community edition is single-tenant, which is fine for local
          development.
    \item Redis uses \ic{--appendonly yes} for crash-safe persistence without
          full RDB snapshots.
\end{itemize}

\begin{warning}
Never bind ports to \ic{0.0.0.0} (the default when you omit the IP).
Doing so exposes your databases---including their credentials---to every device
on your network.  The \ic{127.0.0.1} prefix is a deliberate security measure.
\end{warning}


% ─────────────────────────────────────────────────────────────────────
\section{Volume Persistence}
\label{sec:volume-persistence}

Docker Compose creates four \emph{named volumes} at the bottom of the file:
\ic{postgres\_data}, \ic{neo4j\_data}, \ic{neo4j\_logs}, and \ic{redis\_data}.
Named volumes survive container restarts and even container recreation.

There are two ways to stop the services:

\begin{lstlisting}[style=bash, numbers=none]
# Stop containers, KEEP all data
docker-compose down

# Stop containers and DESTROY all volumes
docker-compose down -v
\end{lstlisting}

The first command is safe to run daily---your databases, user accounts,
migrations, and graph data are all preserved.

\begin{warning}
\texttt{docker-compose down -v} \textbf{wipes everything}: PostgreSQL rows,
Neo4j nodes and edges, Redis caches, and the Alembic migration stamp.  After
running it, you must re-register users, re-run \ic{alembic upgrade head}, and
re-import any data.  Only use \ic{-v} when you intentionally want a clean
slate.  The same warning applies to \ic{docker volume prune}.
\end{warning}


% ─────────────────────────────────────────────────────────────────────
\section{Database Migrations}
\label{sec:database-migrations}

\sysname{} uses \textbf{Alembic} to manage PostgreSQL schema changes.  Alembic
is a migration framework built on SQLAlchemy---it generates incremental SQL
scripts from your Python model definitions and applies them in order.

\subsection*{How Alembic is organized}

\begin{forest}
for tree={
    font=\small\ttfamily,
    grow'=0,
    child anchor=west,
    parent anchor=south,
    anchor=west,
    calign=first,
    edge path={
        \noexpand\path [draw, \forestoption{edge}]
        (!u.south west) +(7.5pt,0) |- (.child anchor)\forestoption{edge label};
    },
    before typesetting nodes={
        if n=1 {insert before={[,phantom]}} {}
    },
    fit=band,
    before computing xy={l=15pt},
}
[apps/api/
    [alembic.ini {\normalfont\itshape\color{ContinuumGray} --- database URL, script location}]
    [alembic/
        [env.py {\normalfont\itshape\color{ContinuumGray} --- reads SQLAlchemy models, configures autogenerate}]
        [versions/
            [001\_initial.py]
            [002\_add\_projects.py]
            [\ldots]
        ]
    ]
]
\end{forest}

\begin{itemize}
    \item \codefile{alembic.ini} tells Alembic where to find the database and
          where to find migration scripts.
    \item \codefile{env.py} imports your SQLAlchemy models so that
          \ic{--autogenerate} can diff the Python definitions against the live
          database.
    \item \codefile{versions/} holds numbered Python scripts, each with an
          \ic{upgrade()} and \ic{downgrade()} function.
\end{itemize}

\subsection*{Common commands}

\begin{lstlisting}[style=bash, numbers=none, caption={Alembic commands you will use most.}]
# Apply all pending migrations
pnpm db:migrate          # or: alembic upgrade head

# Create a new migration from model changes
alembic revision --autogenerate -m "add confidence column"

# Roll back the last migration
alembic downgrade -1

# Check current migration state
alembic current
\end{lstlisting}

\subsection*{The stamp head escape hatch}

If the database already has tables (e.g., you ran the app before adding Alembic)
but the \ic{alembic\_version} table is missing or empty, \ic{alembic upgrade head}
will fail with a \ic{DuplicateTableError}.  The fix:

\begin{lstlisting}[style=bash, numbers=none]
alembic stamp head
\end{lstlisting}

This writes the latest revision ID into \ic{alembic\_version} without actually
running any migration scripts---it tells Alembic ``the database is already at
this state, just record it.''

\begin{warning}
Only use \ic{alembic stamp head} when you are certain the live schema matches
the latest migration.  If the schema is actually behind, stamping skips the
missing changes and your application will encounter column-not-found errors at
runtime.  When in doubt, compare the live schema against the migration scripts
manually.
\end{warning}


% ─────────────────────────────────────────────────────────────────────
\section{Redis Key Patterns}
\label{sec:redis-key-patterns}

Redis stores three tiers of cached data, each with a different TTL based on how
quickly the underlying data can change.
Table~\ref{tab:redis-key-patterns} lists every key pattern used by \sysname{}.

\begin{table}[ht]
\centering
\caption{Redis key patterns, TTLs, and purposes.}
\label{tab:redis-key-patterns}
\begin{tabularx}{\textwidth}{l c X}
\toprule
\textbf{Key Format} & \textbf{TTL} & \textbf{Purpose} \\
\midrule
\ic{entity:\{user\_id\}:\{lookup\}:\{name\}} & 5 min &
    Entity resolution cache.  Avoids re-running the 7-stage pipeline
    for repeated mentions.  User-scoped to enforce data isolation. \\
\addlinespace
\ic{llm:\{version\}:extract:\{hash\}} & 24 h &
    LLM response cache.  Keyed by prompt version so that changing the
    extraction prompt automatically invalidates stale entries.
    Global (same input $\rightarrow$ same output). \\
\addlinespace
\ic{emb:\{model\}:\{type\}:\{hash\}} & 30 d &
    Embedding vector cache.  Embedding models are deterministic and
    rarely change, so a long TTL is safe.  Global. \\
\addlinespace
\ic{rate:\{user\_id\}} & 60 s &
    Rate limit counter.  Tracks request count within a sliding window.
    Per-user.  See Chapter~\ref{ch:configuration},
    Section~\ref{sec:rate-limiting}. \\
\bottomrule
\end{tabularx}
\end{table}

\begin{keyinsight}
The entity cache is per-user, but the LLM and embedding caches are global.
This is intentional: entity resolution is user-scoped (different users may have
different entities with the same name), but LLM and embedding outputs are
deterministic functions of their input regardless of who triggered them.
\end{keyinsight}
